\newpage
\subsubsection{Pruebas de integración}

Las pruebas de integración verificarán la correcta comunicación entre
los componentes del sistema, asegurando que los datos fluyan de forma
íntegra y consistente a través de las distintas capas de la
arquitectura. Esta sección presenta el plan de pruebas de integración
que se ejecutará durante la fase de implementación.

\textbf{Herramientas utilizadas:}

\begin{itemize}
    \item \textbf{Micronaut Test:} Framework oficial de Micronaut para
    pruebas de integración del backend. Permite levantar el contexto
    completo de la aplicación en un entorno de pruebas controlado,
    verificando el comportamiento real de los controladores REST,
    casos de uso y repositorios.
    \item \textbf{Testcontainers:} Biblioteca que levanta instancias
    reales de PostgreSQL en contenedores Docker durante la ejecución
    de las pruebas, eliminando la necesidad de mocks de base de datos
    y garantizando que las pruebas reflejen el comportamiento real del
    sistema en producción.
    \item \textbf{JUnit 5:} Framework base para la definición y
    ejecución de los casos de prueba de integración del backend.
    \item \textbf{Newman:} Ejecutor de colecciones Postman desde línea
    de comandos, integrado en el pipeline de GitHub Actions para
    automatizar las pruebas de los endpoints REST en cada despliegue.
\end{itemize}

La Tabla~\ref{tab:plan_pruebas_integracion} presenta los casos de
prueba de integración planificados.

\begin{longtable}{|>{\centering\arraybackslash}m{1.5cm}|>
{\raggedright\arraybackslash}m{3.5cm}|>{\raggedright\arraybackslash}m{4cm}|>
{\raggedright\arraybackslash}m{4cm}|}
\caption{Plan de pruebas de integración del sistema.}
\label{tab:plan_pruebas_integracion} \\
\hline
\textbf{ID} & \textbf{Caso de prueba} & \textbf{Componentes involucrados} & \textbf{Resultado esperado} \\
\hline
\endfirsthead
\multicolumn{4}{c}{{\bfseries Tabla \thetable\ -- continuación}} \\
\hline
\textbf{ID} & \textbf{Caso de prueba} & \textbf{Componentes involucrados} & \textbf{Resultado esperado} \\
\hline
\endhead

PI-01 & Sincronización móvil → API Gateway & App móvil (Retrofit) → API Gateway & La petición POST /sync llega correctamente a API Gateway con el token JWT en el header. \\
\hline
PI-02 & Validación JWT en API Gateway & API Gateway → Amazon Cognito & API Gateway valida el token JWT con Cognito y rechaza peticiones sin token con HTTP 401. \\
\hline
PI-03 & API Gateway → Lambda & API Gateway → AWS Lambda (Micronaut) & La función Lambda se invoca correctamente y procesa la petición en menos de 5 segundos incluyendo cold start. \\
\hline
PI-04 & Lambda → RDS (escritura) & AWS Lambda → Amazon RDS (PostgreSQL) & Los datos de detección se persisten correctamente en PostgreSQL mediante Micronaut Data JPA y Testcontainers. \\
\hline
PI-05 & Lambda → S3 (carga de imagen) & AWS Lambda → Amazon S3 & La imagen se carga correctamente en S3 y se retorna la URL pública para su almacenamiento en RDS. \\
\hline
PI-06 & Lambda → RDS (lectura) & AWS Lambda → Amazon RDS & La consulta del historial de detecciones retorna los registros correctos desde PostgreSQL. \\
\hline
PI-07 & SageMaker → S3 (publicación modelo) & Amazon SageMaker → Amazon S3 & El modelo \texttt{.tflite} se publica correctamente en S3 y es accesible para descarga desde el dispositivo móvil. \\
\hline
PI-08 & App móvil → S3 (descarga modelo) & App móvil → Amazon S3 & El dispositivo descarga correctamente la nueva versión del modelo \texttt{.tflite} y lo reemplaza localmente. \\
\hline
PI-09 & Angular → API Gateway (petición protegida) & Angular (HttpInterceptor) → API Gateway & El interceptor adjunta el JWT correctamente y la petición recibe respuesta HTTP 200 con datos. \\
\hline
PI-10 & CloudFront → S3 (frontend) & Amazon CloudFront → Amazon S3 & CloudFront sirve los artefactos estáticos de Angular con HTTPS y caché correctamente configurado. \\
\hline
PI-11 & Flujo completo de sincronización & App móvil → API Gateway → Lambda → RDS → S3 & El ciclo completo de sincronización se ejecuta sin errores y el estado en SQLite cambia a \textbf{Sincronizado}. \\
\hline
PI-12 & Flujo completo de consulta web & Angular → API Gateway → Lambda → RDS & La plataforma web consulta y muestra correctamente el historial de detecciones sincronizadas. \\
\hline
PI-13 & TFLite + OpenCV en dispositivo & TFLite (EfficientNet-B3) → OpenCV (análisis cromático) & El modelo TFLite y el análisis cromático de OpenCV se ejecutan de forma concurrente y retornan resultados en menos de 3 segundos sin conexión a internet. \\
\hline
PI-14 & SQLite → Room (escritura local) & Room → SQLite & Los registros de detección se persisten correctamente en SQLite mediante Room, incluyendo estado, coordenadas, timestamp y referencia de imagen. \\
\hline
PI-15 & SQLite → Room (lectura y filtrado) & Room → SQLite & Las consultas de registros pendientes de sincronización retornan únicamente los registros con estado \textbf{Pendiente} o \textbf{Error de sincronización}. \\
\hline
PI-16 & Token refresh en Angular & Angular (HttpInterceptor) → Amazon Cognito & Cuando el Access Token expira, el HttpInterceptor solicita automáticamente un nuevo token a Cognito usando el Refresh Token sin interrumpir la sesión del usuario. \\
\hline
PI-17 & D3.js con datos reales de la API & Angular (D3.js) → API Gateway → Lambda → RDS & Las visualizaciones de D3.js renderizan correctamente los datos de distribución de focos de infestación y evolución temporal obtenidos desde la API. \\
\hline
\end{longtable}